\section{Displacement and Stress Response}

\subsection{Selection of the Point or Area of Study}

% Justify the choice of the study point or area based on the results
% of the modal analysis and the excitation defined in the harmonic analysis.
% Include a screenshot identifying its location.

% TODO: Include screenshot of the selected point
% \begin{figure}[H]
%     \centering
%     \includegraphics[width=0.8\textwidth]{study_point}
%     \caption{Location of the selected point or area of study.}
%     \label{fig:study_point}
% \end{figure}

\subsection{Displacement Response}

% Present a displacement vs.\ frequency graph.
% The graph must include:
% - Horizontal axis with frequency in Hz
% - Vertical axis with displacement in consistent units
% - Identification of the main resonance peaks
% - Technical commentary on the results

% TODO: Include displacement vs.\ frequency graph
% \begin{figure}[H]
%     \centering
%     \includegraphics[width=0.85\textwidth]{displacement_frequency}
%     \caption{Displacement response as a function of frequency.}
%     \label{fig:displacement_frequency}
% \end{figure}

\subsection{Stress Response}

% Present a stress (e.g.\ von Mises) vs.\ frequency graph.
% Identify the main peaks and relate them to the natural frequencies.

% TODO: Include stress vs.\ frequency graph
% \begin{figure}[H]
%     \centering
%     \includegraphics[width=0.85\textwidth]{stress_frequency}
%     \caption{Equivalent stress response as a function of frequency.}
%     \label{fig:stress_frequency}
% \end{figure}

% Relate the observed peaks to the natural frequencies obtained in the modal analysis.